\author{OTS}
\documentclass[fontsize=14pt]{mybib}

\newcommand{\Name}{Fredy}
\graphicspath{{../../../assets/images/}}
\setheaderlogo{mnr-green.png}

% ensure scrlayer-scrpage has sufficient footheight
\setlength{\footheight}{30.4pt}

\begin{document}
\begin{bibelbox}{SCHL}{Kol}{3:15-17}
    Lasst das Wort des Christus reichlich in euch wohnen in aller Weisheit; lehrt und ermahnt einander und singt mit Psalmen und Lobgesängen und geistlichen Liedern dem Herrn lieblich in eurem Herzen. Und was immer ihr tut in Wort und Werk, das tut alles im Namen des Herrn Jesus und dankt Gott, dem Vater, durch ihn.
\end{bibelbox}
\section{Begrüssung}
Ich möchte auch euch alle recht herzlich zum Gottesdienst begrüssen. Schön, dass ihr gekommen seid, um unserem \herr N zu danken, ihn zu loben und zu preisen. Nick und Erika lassen sich entschuldigen, dass sie heute nicht hier sein können. Nick musste kurzfristig nach Ungarn fahren, weil eine Freundin von ihm verstorben ist. 

Auch dich \Name{} will ich herzlich begrüssen. 
\begin{itemize}
    \item[] \beten[b, p]{Beten zur Begrüssung}
    \item[] und anschliessend singen wir zusammen das Lied
    \item[] \lied[b]{Du bist das Licht der Welt.}
\end{itemize}

\section{Informationen}
\begin{itemize}
    \item \info[b, p]{Bibel und Gebetsabend:} Do., 11. Juni 2026 20:00 Uhr Bibel und Gebetsabend mit Samuel Rindlisbacher zu Markus 8,10-13.\\
    Vom 22.6. bis 19.7.2026 sind Janina und ich auf der Alp Erner Galen zum Schafe hüten. Die Gottesdienste während der Zeit werden normal weiter geführt, aber am Donnerstagabend, findet nur ein Gebetsabend statt ohne Liveübertragung. Wenn jemand kommt, der die Liveübertragung auf den Beamer projizieren kann, dann findet die Übertragung statt. Wenn niemand da ist, der weiss wie das funktioniert, findet nur das Gebet statt. Marie-Therese wird ihr Tablet mitnehmen, auf dem könnt ihr dann die Übertragung aus Zürich anschauen. Janina und ich werden die Übertragung auf 2400 über Meer anschauen.
    \item \info[b, p]{Nächster Gottesdienst:} So, 14. Juni 2026 Livestream aus Dübendorf
    \item \info[b, p]{Kollekte:} Die Kollekte wird hinten in der grünen Box gesammelt und wird vollumfänglich für den Bau dieser Gemeinde eingesetzt.
\end{itemize}

\section{Input}
\begin{spacing}{1.5}
   \begin{block}[Gottesnamen]
    Machen wir weiter mit den Namen Gottes. Wir haben die drei Hauptnamen YHWH (Jahwe oder Jehova), Adonai und Elohim. In der Bibel finden wir aus diesen Hauptnamen diverse zusammengesetzte Namen, die oft die Eigenschaften unseres grossen Gottes beschreiben. Zwei haben wir uns schon angeschaut, Jahwe Zebaoth und Jahwe Adonai. Diese Namen haben wir angeschaut, weil sie so häufig in der Bibel auftauche. Den Namen heute wollen wir anschauen, nicht weil er so oft auftaucht, sondern weil er so häufig missbraucht wird. 

    Heute wollen wir uns Jahwe Rapha \enquote{der \herr{} dein Arzt} anschauen.  
    \end{block}
    \begin{block}[Der Herr dein Arzt]
        Dieser Name kommt nicht oft in der Bibel vor, aber er wird oft gebraucht und auch missbraucht. Als Eigennamen kommt er so nur in \bibleverse{IIMos}(15:26) vor:
        \begin{bibelbox}{SCHL}{IIMos}{15:26}
            und er sprach: Wenn du der Stimme des \herr N, deines Gottes, eifrig gehorchen wirst und tust, was ihm recht ist, und seine Gebote zu Ohren fasst und alle seine Satzungen hältst, so will ich keine der Krankheiten auf dich legen, die ich auf Ägypten gelegt habe; denn ich bin der \herr{} dein Arzt.
        \end{bibelbox}
        
        Dieser Vers wird so oft zitiert um zu sagen, dass man nur die Gebote von Gott halten muss und dann wird Gott einem alle Krankheiten heilen. Oder anders gesagt, \enquote{wenn du nur genügend glaubst, wirst du nicht krank}, also \enquote{wenn du krank wirst, glaubst du zu wenig}. Es wird nicht so direkt gesagt, aber wenn du diesen Bibelvers einer kranken Person vorliest, dann wird dieser von dieser Person oft auch so verstanden. Fragt mal Fredy oder Nick, was es mit einem selber macht, wenn man diesen Vers hingeworfen bekommt.

        Ich bin $100\%$ überzeugt, dass Gott heilt, dass Gott geheilt hat und das Gott immer noch heilt. Wieso auch nicht? Er hat ja uns geschaffen, so kann er uns auch reparieren. 
        
        Man liest immer wieder Geschichten, dass Personen geheilt wurden. Aber ich habe noch nie gelesen, oder gehört, dass Gott nur geheilt hat, weil die Person alle Gebote gehalten hätte. Auch im alten Testament nicht. Zum Beispiel Naeman in \bibleverse{IKoen}(5:), er wurde von Gott geheilt, obwohl er nicht mal Jude war. Und im Neuen Testament wird dann berichtet, dass dieser der einzige Mensch in Israel war, der je vom Aussatz geheilt wurde. Jesus sagt in:
        \begin{bibelbox}{SCHL}{Lk}{4:27}
            und viele Aussätzige waren in Israel zur Zeit des Propheten Elisa; aber keiner von ihnen wurde gereinigt, sondern nur Naeman, der Syrer.
        \end{bibelbox}

        Was auch noch bemerkenswert ist: Wenn Gott heilt, wurden diese Menschen immer vollständig wieder hergestellt. Egal ob 40 Jahre gelähmt oder Tod, sie standen auf und konnten direkt wieder umherspazieren.

        \betonung[b, p]{Warum heilt Gott Menschen? } Gott heilt Menschen zu seiner Ehre, es geschieht, dass Gott verherrlicht wird. Es geht nicht darum, in Talkshows aufzutreten oder mit Bücher Geld zu verdienen. Viele Geschichten, die im Umlauf sind, sollte man immer mit einer gewissen Skepsis anschauen. Auch bei sogenannten Nahtoderfahrungen. Am besten ist, man freut sich und betet für diese Personen. Sie wissen am besten was und wie ihnen passiert ist. \betonung[b, p]{Verherrlichen Sie damit Jesus nicht, dann nehmt Abstand davon. }

        Wir brauchen auch keine Heilungsveranstaltungen, wo Wunderheiler einem die Hand auflegen. Solche Heilungen sind nie von Gott. Es gibt auch keine medizinisch bezeugte Heilung aus solchen Veranstaltungen. Aber was es Massenhaft gibt, sind Heiler die als Gauner überführt wurden.

        Ach ja, auch Ärzte können uns nicht heilen. Ärzte können uns Medikamente verschreiben, oder unseren Körper so versorgen, dass dieser sich selbst wieder heilen kann. Unheilbare Krankheit heisst, an einer Krankheit zu leiden, die der Körper nicht heilen kann. Wir bekommen dann Medikamente oder Therapien, die helfen mit dieser Krankheit zu leben. Auch das ist von Gott geschenkt und wir dürfen uns sollen diese Möglichkeiten auch nutzen. 

        Zurück zum Namen. Der Name verspricht in \bibleverse{IIMos}(15:26), dass Gott uns heilen wird, wenn wir seine Gebote halten. Im Alten Testament hat das irdische Schicksal des Volkes Israel am Halten der Gebote gehangen. Der sogenannte Mosebund, wir haben über diesen Bund im Winter gesprochen. Er ist der einzige Bund von Gott mit Menschen, der auf Bedingungen ihrer seit's beruht. Wir sind jetzt im neuen Bund, in dem uns der Herr versprochen hat, uns ein neues Herz zu geben. 

        Jesus kam auf die Welt und heilte viele Menschen. Um zu zeigen, dass er der Messias ist und so die Schrift erfüllt wird. So lässt er Johannes dem Täufer ausrichten,
        \begin{bibelbox}{SCHL}{Matt}{11:5}
        Blinde werden sehend und Lahme gehend Aussätzige werden rein und Taube hören, Tote werden auferweckt, und Armen wird das Evangelium verkündigt.
        \end{bibelbox}

        War das, das Ziel Jesus, Menschen von ihren Krankheiten zu heilen? Bedeutet \betonung[b, p]{das} \enquote{der Herr dein Arzt}. Nein! Der Herr dein Arzt heilt dich, in dem er sich selber schinden liess und für uns am Kreuz gestorben ist. Wenn es im alten Testament heisst, ich werde euch segnen, wenn ihr euch an meine Gebote haltet, galt das nur für irdische Segnungen. Aber im Alten wie im Neuen Testament, werden wir durch den Glauben für ewig geheilt.

        Um geheilt zu werden, müssen wir auf das Kreuz blicken, auf den der für unsere Sünden gestorben ist. Er ist unser Arzt, der uns heilt. Nein, nicht von irdischen Krankheiten, sondern von der Krankheit der Sünde. Diese Krankheit ist wie Aussatz in der Ewigkeit und so aussätzig können wir nicht vor Gott treten. Für alle die hier auf Erden sich von ihren liebsten verabschieden mussten, wenn diese Gott als ihren Arzt von Herzen angenommen haben, sind sie für die Ewigkeit von diesem Aussatz der Sünde geheilt und dürfen so vor das Angesicht Gottes durch Jesus Christus treten. Das ist unser grösster Trost und unsere wichtigste Heilung. So steht in Psalm \bibleverse{Ps}(103:1-3):
        \begin{bibelbox}{SCHL}{Ps}{103:1-3}
            Von David. Lobe den Herrn, meine Seele, und alles, was in mir ist, seinen heiligen Namen! Lobe den Herrn, meine Seele, und vergiss nicht, was er dir Gutes getan hat! Der dir alle deine Sünden vergibt und heilt alle deine Gebrechen;
        \end{bibelbox}
        
        Der aktivste Fitness Guru oder gesündeste Mensch der diesen Arzt nicht annimmt, wird in der Ewigkeit mit dem Aussatz der Sünde behaftet sein und so für ewig von Gott getrennt bleiben.

        Lasst Gott euer Arzt sein. Dieser Arzt ist für unsere Heilung am Kreuz gestorben. Tragt das in die Welt hinaus. Denn wir haben keinen Fachkräftemangel an Heilern, weil Jesus ist für \betonung[b, p]{jeden} Menschen gestorben. 
        \begin{bibelbox}{SCHL}{Joh}{3:16}
        Denn so sehr hat Gott die Welt geliebt, dass er seinen eingeborenen Sohn gab, damit \betonung[b, p]{jeder}, der an ihn glaubt, nicht verloren geht, sondern ewiges Leben hat.
        \end{bibelbox}
        
        Wenn ein Arzt dir ein Rezept ausstellt, musst du dieses annehmen. So musst du dir selber die gleiche Frage stellen, die Jesus auch an Martha stellte:
        \begin{bibelbox}{SCHL}{Joh}{11:25}
            Jesus spricht zu ihr: Ich bin die Auferstehung und das Leben. Wer an mich glaubt, wird leben, auch wenn er stirbt; und jeder, der lebt und an mich glaubt, wird in Ewigkeit nicht sterben. Glaubst du das?
        \end{bibelbox}
        \betonung[b, p]{Was ist deine Antwort?} Nur wenn du diese mit Ja beantwortest, wirst du Heilung erfahren.

        Das ist unser Arzt. \betonung[b, p]{\herr{} dein Arzt} Und wenn wir heute an den Tisch des Herrn treten, denken wir daran, wie der Herr unser Arzt uns von unseren Sünden geheilt hat, damit wir für die Ewigkeit vor Gott treten können.

        Lesen wir noch die Verse \bibleverse{Ps}(147:1-3)
        \begin{bibelbox}{SCHL}{Ps}{147:1-3}
        Lobt den Herrn! Denn es ist gut, unserem Gott zu lobsingen: Es ist lieblich, es gebührt ihm Lobgesang. Der Herr baut Jerusalem; die Zerstreuten Israels wird er sammeln. Er heilt, die zerbrochenen Herzens sind, und verbindet ihre Wunden.
        \end{bibelbox}
        Darauf dürfen wir uns freuen, wenn der Herr wieder kommt, oder wenn er uns zu sich ruft! Amen
    \end{block}
    
    \lied[b]{Denn ich bin gewiss}.
\end{spacing}
\section{Predigt}
\textbf{Nach der Predigt}\newline
Danken für die Predigt.
\section{Abendmahl}
\begin{itemize}
    \item [] Beten für das Brot
    \item [] Beten für den Wein
\end{itemize}
\section{Abschluss}
\begin{itemize}
    \item [] \lied[b]{Bald schon kann es sein}.
    \item [] \beten[b]{Zum Abschluss}
\end{itemize}
\begin{bibelbox}{SCHL}{IIKor}{13:13}
    Die Gnade des Herrn Jesus Christus und die Liebe Gottes und die Gemeinschaft des Heiligen Geistes sei mit euch allen! Amen
\end{bibelbox}
\end{document}